% Gerado por docs/gerar_declaracao_ia.py a partir de docs/registro_ia.toml. Não edite à mão.
Neste trabalho foi utilizada a ferramenta de inteligência artificial generativa Claude Code (Claude Opus 5.5, claude-opus-5-5), da Anthropic, em 2 sessões e 2 interações, nas datas 23/11/2026. O \autoref{quadro:uso-ia} resume a finalidade de cada interação e a verificação feita pelo autor, e o \autoref{apendice:registro-ia} reproduz a prompt literal e a saída de cada uma. Todo o conteúdo produzido com auxílio da ferramenta foi revisado pelo autor, que se responsabiliza integralmente pelo trabalho.

\textbf{Sessão S1} -- interface: \pendente{informar a interface usada (ex.: claude.ai/code na web, Claude Code no terminal ou no VS Code) e a versão}. Configurações: \pendente{busca na web, memória/histórico de chats, modo de permissões, plugins}. Instruções de sistema ou de projeto: \pendente{instruções de projeto ou personalizadas que estavam ativas, se houver (o prompt de sistema padrão da ferramenta não é reproduzido)}.

\textbf{Sessão S2} -- interface: Claude Code 2.1.330, extensão do Visual Studio Code, executando no Ubuntu 26.04.1 LTS (WSL 2). Configurações: modo de permissões automático; ferramenta de busca na web disponível, mas não usada; memória persistente do Claude Code ativada (vazia no início da sessão); plugin academico-br habilitado; temperatura e demais parâmetros de amostragem no padrão da ferramenta (não configuráveis pelo usuário). Instruções de sistema ou de projeto: prompt de sistema padrão do Claude Code (não reproduzido). Nenhum CLAUDE.md no projeto no início da sessão; o CLAUDE.md criado nesta sessão passa a valer para as seguintes.

\begin{quadro}[htbp]
    \caption{Interações com a ferramenta de inteligência artificial}
    \label{quadro:uso-ia}
    \centering
    \footnotesize
    \begin{tabular}{|c|c|c|p{4.3cm}|p{5.8cm}|}
        \hline
        \textbf{Nº} & \textbf{Sessão} & \textbf{Data} & \textbf{Finalidade} & \textbf{Verificação humana} \\ \hline
        I1 & S1 & 23/11/2026 & Geração do código: estrutura inicial do repositório (árvores AVL e Rubro-Negra em C, benchmark, validação e análise em Python) & \pendente{descrever o que foi conferido (ex.: leitura do código, execução de ./main validar, comparação com o CLRS), corrigido ou descartado.} \\ \hline
        I2 & S2 & 23/11/2026 & Estrutura do relatório em LaTeX (abnTeX2) e criação deste registro de uso de IA & \pendente{descrever o que foi conferido, corrigido ou descartado.} \\ \hline
    \end{tabular}
    \fonte{Elaborado pelo autor (2026).}
\end{quadro}
